% ch00-primer
\section{من الصفر: يعني إيه نظام برمجيّ؟}

أهلًا بك في أوّل فصل من الكتاب. لو أنت لا تعرف أيّ شيء عن البرمجة، فأنت في المكان
الصحيح تمامًا. هدفنا هنا بسيط: أن تُغلق هذا الفصل وأنت فاهم \emph{إيه هو} نظام
\en{ScholarPath}، ومن أيّ قطع كبيرة يتكوّن، وكيف تتكلّم هذه القطع مع بعضها. مش هنغرق
في تفاصيل تقنية عميقة دلوقتي؛ ده مجرّد تمهيد نبني عليه باقي الفصول. خُد نفَسًا، وتعال
نمشي \emph{واحدة واحدة}.

\subsection{إيه هو \en{ScholarPath} أصلًا؟}

تخيّل طالبًا مصريًّا في سنته الأخيرة بالجامعة، عايز يكمّل دراسته بمنحة في الخارج. هو
يعرف أنّ المنح موجودة، لكنّه \emph{تايه}: المنح متناثرة على عشرات المواقع بلغات
مختلفة، ومواعيد التقديم متضاربة، والأوراق المطلوبة غامضة، ولا يوجد من يرشده هل هو
مؤهَّل أصلًا لهذه المنحة أم لا. النتيجة المعتادة: يضيع الميعاد، أو يقدّم بشكل ناقص،
أو ييأس ويستسلم.

\begin{idea}[الفكرة: ما هو \en{ScholarPath}؟]
\en{ScholarPath} هو \textbf{منصّة إلكترونية} (موقع ويب) تجمع رحلة المنحة الدراسية
كلّها في مكان واحد. المنصّة \textbf{تكتشف} المنح المناسبة للطالب، و\textbf{تتابع} خطوات
تقديمه عليها، وتتيح له \textbf{حجز جلسة فيديو} مع استشاريّ متخصّص (جلسة مدفوعة)، وفيها
\textbf{مجتمع} يسأل فيه ويتناقش، و\textbf{مساعد ذكاء اصطناعيّ} يجاوب على أسئلته.
والمنصّة \textbf{ثنائية اللغة}: تعمل بالعربية (من اليمين لليسار) وبالإنجليزية.
\end{idea}

المنصّة تخدم \textbf{أربعة أنواع من المستخدمين}، ونسمّي كلّ نوع \en{«role»} أي «دور».
كلّ دور يرى شاشات مختلفة وله صلاحيّات مختلفة:

\begin{center}
\begin{tabular}{L{3.2cm} L{9cm}}
\toprule
\textbf{\color{spPrimary}الدور} & \textbf{\color{spPrimary}بيعمل إيه على المنصّة} \\
\midrule
الطالب \en{(Student)} & يبحث عن المنح، يتابع تقديماته، يحجز جلسات استشاريّ، يشارك في المجتمع، ويسأل المساعد الذكيّ. \\
الاستشاريّ \en{(Consultant)} & خبير يعرض أوقاته المتاحة، يستقبل الحجوزات، ويقدّم جلسات فيديو مدفوعة للطلاب. \\
الجهة المانحة \en{(Scholarship Provider)} & تنشر منحها على المنصّة، وتراجع طلبات الطلاب الخاصّة بها. \\
الأدمن \en{(Admin)} & المشرف الذي يوافق على المحتوى، يدير المستخدمين، ويراقب سلامة المنصّة. \\
\bottomrule
\end{tabular}
\end{center}

\begin{keyp}[نقطة مهمة: حجم النظام]
\en{ScholarPath} ليس تطبيقًا صغيرًا. هو مقسَّم إلى \textbf{18 وحدة كبرى} (نسمّيها
\en{modules} أو \en{epics}) مرقّمة من \en{PB-001} إلى \en{PB-018}، تغطّي الرحلة كاملة
من إنشاء الحساب حتى تحليلات البيانات. وبلغة الأرقام: 211 \textbf{متطلَّبًا وظيفيًّا}
(\en{functional requirement})، و203 \textbf{حالة استخدام} (\en{use case})،
و159 \textbf{قصّة مستخدم} (\en{user story}). لا تقلق من هذه المصطلحات الآن — هي مجرّد
طرق منظَّمة لكتابة «النظام لازم يعمل كذا».
\end{keyp}

\subsection{يعني إيه برنامج ويب؟ الصالة والمطبخ والمخزن}

أيّ موقع ويب تتعامل معه (فيسبوك، بنكك، \en{ScholarPath}) يتكوّن في الحقيقة من
\textbf{ثلاث قطع كبيرة} تتعاون. أسهل طريقة لفهمها هي تشبيه \textbf{المطعم}:

\begin{center}
\begin{tabular}{L{3.4cm} L{4.3cm} L{4.6cm}}
\toprule
\textbf{\color{spPrimary}في المطعم} & \textbf{\color{spPrimary}في النظام} & \textbf{\color{spPrimary}وظيفته} \\
\midrule
الصالة والجرسون & الواجهة \en{(Frontend)} & ما تراه وتلمسه في المتصفّح: الأزرار، الشاشات، النماذج. \\
المطبخ & الخلفية \en{(Backend)} & يستقبل «الطلبات»، ينفّذ قواعد العمل، ويجهّز النتيجة. \\
المخزن والثلّاجة & قاعدة البيانات \en{(Database)} & تخزّن كلّ شيء بشكل دائم: المستخدمون، المنح، الحجوزات. \\
\bottomrule
\end{tabular}
\end{center}

\begin{idea}[الفكرة: كيف تتعاون القطع الثلاث؟]
عندما تطلب أكلة، أنت لا تدخل المطبخ ولا المخزن بنفسك. تقول للجرسون (الواجهة) ما تريد،
الجرسون ينقل الطلب للمطبخ (الخلفية)، المطبخ يأخذ المكوّنات من المخزن (قاعدة البيانات)،
يطبخها، ثمّ يعيد الطبق للجرسون الذي يضعه أمامك. \textbf{الواجهة} تعيش داخل متصفّحك على
جهازك، بينما \textbf{الخلفية} و\textbf{قاعدة البيانات} تعيشان على حاسوب بعيد قويّ اسمه
\en{«server»} (الخادم).
\end{idea}

\begin{note}
لماذا نفصل هذا الفصل؟ لأنّ لكلّ قطعة أدوات مختلفة. واجهة \en{ScholarPath} مكتوبة بلغة
\en{React} وتعمل في المتصفّح، وخلفيّته مكتوبة بلغة \en{C\#} على منصّة \en{.NET}،
وقاعدة بياناته هي \en{SQL Server}. سنشرح كلّ اسم من هذه بعد قليل.
\end{note}

\subsection{يعني إيه \en{API}؟ ورحلة ضغطة زرّ واحدة}

طيب، بأيّ «لغة» يتكلّم الجرسون مع المطبخ؟ في عالم الويب، الواجهة والخلفية يتكلّمان
عبر ما يُسمّى \en{API}.

\begin{idea}[الفكرة: ما هو \en{API}؟]
\en{API} اختصار لـ \en{Application Programming Interface}، أي \textbf{«واجهة برمجية»}.
تخيّلها \textbf{قائمة الطعام} في المطعم: قائمة محدَّدة من الطلبات المسموح بها، ولكلّ طلب
شكل معروف. الواجهة لا تكلّم قاعدة البيانات مباشرة أبدًا؛ هي فقط تطلب من الخلفية عبر هذه
القائمة. مثلًا: «هات لي المنح صفحة رقم 1»، أو «احجز الجلسة الفلانية».
\end{idea}

وكلّ طلب-وردّ بينهما يمشي على بروتوكول اسمه \en{HTTP}. الطلب اسمه \en{request}،
والردّ اسمه \en{response}. لِنتتبَّع معًا \textbf{رحلة ضغطة زرّ واحدة}: الطالب فتح صفحة
المنح وضغط «بحث».

\begin{enumerate}
\item \textbf{المتصفّح (الواجهة)} يجهّز \en{request} ويرسله للخلفية: «\en{GET} المنح».
\item \textbf{الخلفية} تستقبل الطلب، وتتأكّد أنّ المستخدم مسجَّل الدخول (صلاحيّات).
\item الخلفية تسأل \textbf{قاعدة البيانات}: «هات المنح المنشورة، صفحة 1».
\item قاعدة البيانات ترجع الصفوف، والخلفية تحوّلها إلى \en{response} بصيغة \en{JSON}.
\item المتصفّح يستقبل الردّ ويرسم بطاقات المنح على الشاشة. كلّ ذلك في أقلّ من ثانية.
\end{enumerate}

\begin{note}[ما هو \en{JSON}؟]
\en{JSON} هي طريقة قياسية لكتابة البيانات كنصّ منظَّم بين أقواس، يفهمها الطرفان بسهولة.
هي «لغة التغليف» التي يسافر بها الردّ من الخلفية إلى الواجهة.
\end{note}

خلّينا نشوف الشكل التوضيحيّ للطلب والردّ. \textbf{المقتطف التالي توضيحيّ فقط} (مبسَّط
لنفهم الشكل، وليس منسوخًا من ملفّ):

\begin{english}
\begin{lstlisting}[language={}]
// Request — the browser asks the server for scholarships
GET /api/scholarships?page=1&fundedOnly=true
Accept-Language: ar

// Response — the server replies (200 OK) with JSON
{
  "items": [
    { "id": "a1b2-3c4d", "title": "Chevening", "fundingType": "FullyFunded" }
  ],
  "totalCount": 42
}
\end{lstlisting}
\end{english}

هنشوف إيه في المقتطف ده؟ السطر الأوّل \code{GET /api/scholarships} هو \emph{الطلب}: كلمة
\en{GET} تعني «هات لي» (قراءة بيانات)، والباقي هو العنوان مع فلاتر البحث. السطر
\code{Accept-Language: ar} يقول للخلفية «أنا عايز الردّ بالعربي». ثمّ يأتي \emph{الردّ}:
الرقم \code{200} معناه «تمام، نجح الطلب»، وبعده كتلة \en{JSON} فيها قائمة \code{items}
(المنح) وعددها الكلّي \code{totalCount}.

الآن نرى \textbf{الكود الحقيقيّ} من المشروع الذي يقوم بهذا الطلب. أوّلًا من جهة
\textbf{الواجهة}: هذا المقتطف من ملفّ الخدمة \code{client/src/services/api/scholarships.ts}
الذي يرسل الطلب فعلًا:

\begin{english}
\begin{lstlisting}[language=Java]
const { data } = await apiClient.get<PaginatedListWireDto<ScholarshipWireDto>>(
  "/api/scholarships",
  { params },
);
return {
  items: data.items.map(toListItem),
  total: data.totalCount,
  totalPages: data.totalPages,
};
\end{lstlisting}
\end{english}

هنشوف إيه هنا؟ الدالة \code{apiClient.get} هي «الجرسون» الذي يذهب للخلفية على العنوان
\code{"/api/scholarships"} ومعه فلاتر البحث \code{params}. كلمة \code{await} معناها «استنّى
الردّ يرجع». عندما يعود الردّ نأخذ منه \code{data.items} (قائمة المنح) ونجهّزها لعرضها.

وفي \textbf{الجهة المقابلة} على الخلفية، هذا هو الجزء الذي \emph{يستقبل} الطلب، من ملفّ
\code{ScholarshipsController.cs}:

\begin{english}
\begin{lstlisting}[language={[Sharp]C}]
[HttpGet]
public async Task<ActionResult<PaginatedList<ScholarshipDto>>> Get(
    [FromQuery] GetScholarshipsQuery query)
{
    // read the requested language from the Accept-Language header
    var headerLang = Request.Headers["Accept-Language"].ToString()
        .Split(',').FirstOrDefault() ?? "en";
    var lang = headerLang.StartsWith("ar",
        StringComparison.OrdinalIgnoreCase) ? "ar" : "en";

    var updatedQuery = query with { Language = lang };
    return await mediator.Send(updatedQuery);
}
\end{lstlisting}
\end{english}

هنشوف إيه هنا؟ السطر \code{[HttpGet]} يقول «هذه الدالة تردّ على طلبات \en{GET}». الدالة
اسمها \code{Get} وتستقبل شروط البحث في \code{query}. السطرين في الوسط يقرآن ترويسة
\code{Accept-Language} ليعرفا هل يردّ بالعربي أم الإنجليزي. أخيرًا \code{mediator.Send}
يمرّر الطلب للطبقة الداخلية التي تجلب البيانات وترجعها. لاحظ التطابق الجميل: الواجهة
نادت \code{"/api/scholarships"}، وهنا \code{[HttpGet]} على نفس العنوان يردّ عليها.

\begin{secnote}
لاحظ سطر \code{Request.Headers} وقبله التحقّق من تسجيل الدخول (الصفة \en{[Authorize]} على
هذا المتحكِّم). في \en{ScholarPath} \emph{لا يُقدَّم أيّ شيء دون التأكّد من هوية المستخدم
وصلاحيّته} — قاعدة أمان أساسية اسمها «البوّابة المغلقة» \en{(gated)}.
\end{secnote}

\subsection{يعني إيه «طبقة» \en{(layer)}؟ ولماذا نقسّم الكود؟}

لاحظت في المقتطف السابق أنّ المتحكِّم لم يجلب البيانات بنفسه، بل «مرّرها للطبقة
الداخلية». هذه فكرة محورية في \en{ScholarPath} اسمها \textbf{الطبقات} \en{(layers)}.

\begin{idea}[الفكرة: لماذا طبقات؟]
بدل أن نكتب كلّ شيء (شكل الشاشة + قواعد العمل + كلام قاعدة البيانات) في مكان واحد
متشابك، نقسّم الكود إلى \textbf{طبقات} لكلٍّ منها مسؤولية واحدة واضحة. تخيّل المطعم مرة
أخرى: الجرسون لا يطبخ، والطبّاخ لا يزرع الخضار. هذا الفصل يجعل الكود \textbf{أسهل في
الفهم، والاختبار، والتعديل} دون كسر باقي أجزاء النظام.
\end{idea}

خلفيّة \en{ScholarPath} مقسَّمة إلى \textbf{أربع طبقات}، والقلب فيها هو طبقة \code{Domain}.
هذا مقتطف حقيقيّ من قلب النظام، من ملفّ \code{ScholarPath.Domain/Entities/Scholarships.cs}،
يمثّل «المنحة» ككائن أعمال خالص:

\begin{english}
\begin{lstlisting}[language={[Sharp]C}]
public class Scholarship : AuditableEntity, ISoftDeletable
{
    public string TitleEn { get; set; } = default!;
    public string TitleAr { get; set; } = default!;   // bilingual title
    public ScholarshipStatus Status { get; set; } = ScholarshipStatus.Draft;
    public DateTimeOffset Deadline { get; set; }
    public FundingType FundingType { get; set; }
    public decimal? FundingAmountUsd { get; set; }
    // ... more properties omitted
}
\end{lstlisting}
\end{english}

هنشوف إيه هنا؟ هذا الكائن \code{Scholarship} يصف «ما هي المنحة»: لها عنوان بالعربي
\code{TitleAr} وبالإنجليزي \code{TitleEn}، وحالة \code{Status} (مسودّة، منشورة...)، وميعاد
نهائيّ \code{Deadline}، ونوع تمويل \code{FundingType}. المهمّ: هذا الكائن \textbf{لا يعرف
أيّ شيء} عن قاعدة البيانات ولا عن الشاشات — هو قواعد عمل خالصة. سنتعمّق في هذا التصميم
(اسمه \en{Clean Architecture}) في الفصل التالي؛ يكفيك الآن أن تعرف \emph{أنّ} الكود
مقسَّم لطبقات، وأنّ \code{Domain} هي القلب.

\subsection{يعني إيه سحابة \en{(cloud)} و\en{Azure}؟}

قلنا إنّ الخلفية وقاعدة البيانات تعيشان على «خادم بعيد». لكن مَن يملك هذا الخادم؟
هنا تأتي فكرة \textbf{السحابة} \en{(cloud)}.

\begin{idea}[الفكرة: ما هي السحابة؟]
«السحابة» ببساطة تعني \textbf{حواسيب شركة كبيرة نستأجرها عبر الإنترنت}، بدل أن نشتري
خوادم ونضعها في غرفة عندنا. ندفع مقابل ما نستخدمه فقط، وهم يتولّون الكهرباء والتبريد
والصيانة. \en{ScholarPath} يعيش على سحابة \textbf{مايكروسوفت} اسمها \en{Azure}.
\end{idea}

نستخدم من \en{Azure} عدّة خدمات، كلمة واحدة لكلٍّ منها: \en{App Service} (يشغّل
الخلفية)، \en{SQL Database} (قاعدة البيانات)، \en{Blob Storage} (تخزين الملفّات
المرفوعة)، \en{Azure OpenAI} (محرّك الذكاء الاصطناعيّ)، و\en{Key Vault} (خزنة كلمات
السرّ). لا داعي لحفظها الآن؛ ستقابلها تباعًا في فصول لاحقة.

\subsection{مكوّنات المنصّة: مَن يعمل ماذا}

جمعنا في الجدول التالي \textbf{أهمّ التقنيات} التي يقوم عليها \en{ScholarPath}، بسطر
مبسَّط لكلّ واحدة. ستراها مرارًا في الكتاب، فلا بأس أن تعود لهذا الجدول كلّما نسيت:

\begin{center}
\begin{tabular}{L{3cm} L{9.2cm}}
\toprule
\textbf{\color{spPrimary}التقنية} & \textbf{\color{spPrimary}بتعمل إيه} \\
\midrule
\en{.NET 10} & المنصّة (بلغة \en{C\#}) التي تُبنى عليها الخلفية كلّها. \\
\en{React 19} & المكتبة (بلغة \en{TypeScript}) التي تُبنى بها الواجهة في المتصفّح. \\
\en{SQL Server} & قاعدة البيانات التي تخزّن كلّ شيء في جداول. \\
\en{Azure OpenAI} & محرّك الذكاء الاصطناعيّ الذي يشغّل المساعد الذكيّ ويجاوب الأسئلة. \\
\en{Stripe} & بوّابة الدفع الآمنة التي تُحصِّل ثمن جلسات الاستشاريّ بالبطاقة. \\
\en{SignalR} & القناة الفورية للتحديثات اللحظية: الدردشة والإشعارات دون تحديث الصفحة. \\
\en{Hangfire} & مدبّر المهام الخلفية المجدولة (تذكيرات، تنظيف، معالجة مؤجَّلة). \\
\bottomrule
\end{tabular}
\end{center}

\subsection{بنية المشروع على القرص}

كفاية نظريّة، تعال نفتح المشروع فعلًا. الكود مقسَّم إلى جزأين كبيرين على القرص:
مجلَّد \code{server/} للخلفية، ومجلَّد \code{client/} للواجهة. هذه هي شجرة \textbf{الخلفية}
الحقيقية، وفيها الطبقات الأربع التي تكلّمنا عنها:

\begin{english}
\begin{lstlisting}[language={}]
server/src/
|-- ScholarPath.Domain/           # the core: entities + business rules
|   |-- Entities/  Enums/  Events/  Interfaces/  Exceptions/
|-- ScholarPath.Application/       # use cases: Commands / Queries
|   |-- Scholarships/  Applications/  ConsultantBookings/  AI/  ...
|-- ScholarPath.Infrastructure/    # database + external services
|   |-- Persistence/  Migrations/  Services/  Hubs/  Jobs/
`-- ScholarPath.API/               # entry point: Controllers + Middleware
    |-- Controllers/  Middleware/  Program.cs  appsettings.json
\end{lstlisting}
\end{english}

هنشوف إيه هنا؟ القراءة من الداخل للخارج: \code{Domain} هي القلب (كائنات مثل
\code{Scholarship} الذي رأيناه). \code{Application} فيها «حالات الاستخدام» مقسَّمة حسب
الموديول (مجلَّد \code{Scholarships}، مجلَّد \code{ConsultantBookings}...). \code{Infrastructure}
تتكلّم مع قاعدة البيانات والخدمات الخارجية. و\code{API} هي المدخل الذي تصل إليه طلبات
الويب — وفيها مجلَّد \code{Controllers} حيث يعيش المتحكِّم الذي رأيناه.

وهذه شجرة \textbf{الواجهة} (\code{client/src}) باختصار:

\begin{english}
\begin{lstlisting}[language={}]
client/src/
|-- pages/         # one folder per role: student/ consultant/ company/ admin/
|-- components/    # reusable UI pieces (buttons, cards, forms)
|-- services/      # api/ + signalR/  (these talk to the backend)
|-- hooks/  stores/  routes/  locales/  theme/  types/
`-- App.tsx  main.tsx
\end{lstlisting}
\end{english}

هنشوف إيه هنا؟ مجلَّد \code{pages} فيه مجلَّد لكلّ دور (\code{student}، \code{consultant}...).
مجلَّد \code{components} فيه القطع القابلة لإعادة الاستخدام. والأهمّ لموضوعنا: مجلَّد
\code{services} فيه \code{api} — ومنه جاء المقتطف \code{apiClient.get} الذي يكلّم الخلفية.
مجلَّد \code{locales} يحمل ترجمات العربية والإنجليزية.

\begin{qa}
\textbf{سؤال متوقَّع:} ليه الواجهة والخلفية في مجلَّدين منفصلين وبلغتين مختلفتين؟
\newline
\textbf{الإجابة:} لأنّهما يعملان في بيئتين مختلفتين — الواجهة داخل متصفّح المستخدم
والخلفية على الخادم — ولكلٍّ منهما أدواته المثلى. الفصل يسمح لفريق الواجهة وفريق
الخلفية بالعمل بالتوازي، ويسمح باستبدال أيّ جهة دون كسر الأخرى، طالما بقيت لغة الحوار
بينهما (\en{API}) ثابتة.
\end{qa}

\subsection{إزاي تقرأ الكتاب ده — خريطة الـ 3 ساعات}

الكتاب 10 فصول، ومصمَّم بحيث تفهم النظام كاملًا في نحو \textbf{3 ساعات}. لست مضطرًّا
لحفظ أيّ شيء؛ فقط اقرأ بالترتيب، وتوقّف عند الصناديق الملوّنة — فهي تلخّص أهمّ الأفكار.
إليك خريطة تقريبية للوقت:

\begin{center}
\begin{tabular}{L{4.5cm} L{5cm} C{2.4cm}}
\toprule
\textbf{\color{spPrimary}المرحلة} & \textbf{\color{spPrimary}تتعلّم فيها} & \textbf{\color{spPrimary}وقت تقريبيّ} \\
\midrule
هذا الفصل (0) & الصورة الكبيرة والمصطلحات & 20 دقيقة \\
البنية والطبقات & \en{Clean Architecture} والتدفّق & 40 دقيقة \\
الموديولات الأساسية & المنح، التقديم، الحجز، الدفع & 70 دقيقة \\
الميزات المتقدّمة & الذكاء الاصطناعيّ، الفوريّة، التحليلات & 40 دقيقة \\
مراجعة وربط & الأمان والصورة النهائية & 10 دقائق \\
\bottomrule
\end{tabular}
\end{center}

\begin{note}[كيف تستفيد من الصناديق]
الصندوق \textbf{الأزرق} يشرح فكرة، و\textbf{الأخضر} يبرّر قرار تصميم، و\textbf{البنفسجيّ}
نقطة أساسية أو سؤال مناقشة، و\textbf{البرتقاليّ} زاوية أمان. لو كنت مستعجلًا، اقرأ
العناوين والصناديق فقط — تكفي لأخذ فكرة سريعة.
\end{note}

\begin{summary}[الخلاصة في سطور]
\begin{itemize}
\item \en{ScholarPath} منصّة ويب ثنائية اللغة تحلّ تِيه الطالب في عالم المنح: اكتشاف،
متابعة تقديم، حجز جلسات استشاريّ بالفيديو (مدفوعة)، مجتمع، ومساعد ذكاء اصطناعيّ.
\item أربعة أدوار: الطالب، الاستشاريّ، الجهة المانحة، والأدمن — كلٌّ يرى ما يخصّه.
\item أيّ نظام ويب = \textbf{واجهة} في المتصفّح + \textbf{خلفية} على الخادم +
\textbf{قاعدة بيانات} للتخزين (الصالة + المطبخ + المخزن).
\item الواجهة والخلفية يتحاوران عبر \en{API} برسائل \en{HTTP}: \en{request} يخرج،
و\en{response} بصيغة \en{JSON} يعود — كلّ ذلك في أقلّ من ثانية عند ضغطة زرّ.
\item الكود مقسَّم إلى \textbf{طبقات}؛ قلب الخلفية هو طبقة \code{Domain} التي لا تعرف
شيئًا عن الشاشات أو قاعدة البيانات. تفاصيل ذلك في الفصل التالي.
\item يعيش النظام على سحابة \en{Azure}، ويقوم على \en{.NET 10} و\en{React 19}
و\en{SQL Server}، بمساعدة \en{Stripe} و\en{SignalR} و\en{Hangfire} و\en{Azure OpenAI}.
\item على القرص: \code{server/src} فيه الطبقات الأربع، و\code{client/src} فيه الواجهة.
\end{itemize}
\end{summary}
